\chapter{Potenzen und endliche Rechenformeln}
\label{b10guide:ch:calculus}

Alle Zahlen und alle Laufindizes dieses Kapitels liegen in $\mathbb N$.
Endliche Summen und Produkte werden ausschließlich für natürliche
Zahlen definiert. Die Aussagen betreffen ihre arithmetischen Werte.

\section{Potenzen}

\begin{definition}[Natürliche Potenzen]
\label{b10guide:def:powers}
Für $a\in\mathbb N$ sind die Potenzen durch
\[
 a^0=1,\qquad a^{n+1}=a^n a
\]
eindeutig festgelegt. Diese Definition schließt $0^0=1$ ein.
\end{definition}

\begin{theorem}[Potenzgesetze]
\label{b10guide:thm:power-laws}
Für $a,b,m,n\in\mathbb N$ gelten
\[
 a^{m+n}=a^m a^n,\qquad
 (a^m)^n=a^{mn},\qquad
 (ab)^n=a^n b^n.
\]
Außerdem gelten $a^1=a$ und $1^n=1$.
\end{theorem}

\begin{theorem}[Null und eins bei Potenzen]
\label{b10guide:thm:power-special}
Für $a,n\in\mathbb N$ gelten
\[
 a^n=0\ \Longleftrightarrow\ (a=0\ \land\ n>0),
 \qquad
 a^n=1\ \Longleftrightarrow\ (n=0\ \lor\ a=1).
\]
\end{theorem}

\begin{theorem}[Vergleich von Potenzen]
\label{b10guide:thm:power-order}
Für $r>0$ gelten
\[
 a\le b\ \Longleftrightarrow\ a^r\le b^r,
 \qquad
 a<b\ \Longleftrightarrow\ a^r<b^r.
\]
Für $a\ge2$ gelten
\[
 m\le n\ \Longleftrightarrow\ a^m\le a^n,
 \qquad
 m<n\ \Longleftrightarrow\ a^m<a^n,
 \qquad a^n\ge n+1.
\]
Für die Basen $0$ und $1$ gelten die Regeln aus
Satz~\ref{b10guide:thm:power-special}.
\end{theorem}

\section{Endliche Summen und Produkte}

\begin{definition}[Rekursive Summen- und Produktnotation]
\label{b10guide:def:finite-calculus}
Für eine vorgegebene Zuordnung $i\mapsto a_i\in\mathbb N$ setzen wir
\[
 \begin{aligned}
 \sum_{i<0}a_i&=0,&
 \sum_{i<n+1}a_i&=\left(\sum_{i<n}a_i\right)+a_n,\\
 \prod_{i<0}a_i&=1,&
 \prod_{i<n+1}a_i&=\left(\prod_{i<n}a_i\right)a_n.
 \end{aligned}
\]
Nur die angegebenen endlich vielen Werte werden benutzt.
Die Schreibweisen $\sum_{i=0}^{n}a_i$ und $\prod_{i=0}^{n}a_i$
stehen für die Ausdrücke mit $i<n+1$; eine untere Grenze $r$
bezeichnet entsprechend die Werte $a_r,a_{r+1},\ldots$.
\end{definition}

\begin{theorem}[Blockzerlegung und Summenrechnung]
\label{b10guide:thm:finite-calculus-laws}
Für natürliche Werte $a_i,b_i$ und $c,m,n\in\mathbb N$ gelten
\[
 \begin{aligned}
 \sum_{i<m+n}a_i
   &=\sum_{i<m}a_i+\sum_{j<n}a_{m+j},\\
 \prod_{i<m+n}a_i
   &=\left(\prod_{i<m}a_i\right)\left(\prod_{j<n}a_{m+j}\right),\\
 \sum_{i<n}(a_i+b_i)
   &=\sum_{i<n}a_i+\sum_{i<n}b_i,\\
 \sum_{i<n}ca_i&=c\sum_{i<n}a_i,\qquad
 \sum_{i<n}c=nc,\qquad \prod_{i<n}c=c^n.
 \end{aligned}
\]
\end{theorem}

\begin{theorem}[Null- und Positivitätsregeln endlicher Ausdrücke]
\label{b10guide:thm:finite-calculus-zero}
Für natürliche Werte $a_i$ gelten
\[
 \begin{aligned}
 \sum_{i<n}a_i=0
   &\ \Longleftrightarrow\ \forall i<n\ (a_i=0),\\
 \sum_{i<n}a_i>0
   &\ \Longleftrightarrow\ \exists i<n\ (a_i>0),\\
 \prod_{i<n}a_i=0
   &\ \Longleftrightarrow\ \exists i<n\ (a_i=0),\\
 \prod_{i<n}a_i=1
   &\ \Longleftrightarrow\ \forall i<n\ (a_i=1).
 \end{aligned}
\]
Die Aussagen umfassen den leeren Fall $n=0$.
\end{theorem}

\begin{theorem}[Klassische Summenformeln]
\label{b10guide:thm:basic-sums}
Für $n\in\mathbb N$ gelten
\[
 2\sum_{i<n}(i+1)=n(n+1),\qquad
 \sum_{i<n}(2i+1)=n^2,
\]
und
\[
 6\sum_{i<n}(i+1)^2=n(n+1)(2n+1).
\]
Außerdem gilt die Kubensummenformel
\[
 \sum_{i<n}(i+1)^3=\left(\sum_{i<n}(i+1)\right)^2.
\]
Diese Fassungen verwenden nur Addition und Multiplikation in
$\mathbb N$.
\end{theorem}

\begin{theorem}[Endliche geometrische Summe]
\label{b10guide:thm:geometric-sum}
Für $a,n\in\mathbb N$ gilt
\[
 1+a\sum_{i<n}a^i=\sum_{i<n}a^i+a^n.
\]
Für $a\ge1$ kann dieselbe Identität auch als
\[
 (a-1)\sum_{i<n}a^i+1=a^n
\]
geschrieben werden. Die Formel gilt auch für $n=0$ und $a=1$.
\end{theorem}

\section{Fakultäten und Binomialkoeffizienten}

Die Binomialkoeffizienten werden hier rein arithmetisch eingeführt;
als Vergleichsquelle dient \cite{natchoose}.

\begin{definition}[Fakultät]
\label{b10guide:def:factorial}
Die Fakultät wird durch
\[
 0!=1,\qquad (n+1)!=n!(n+1)
\]
definiert. Gleichwertig ist $n!=\prod_{i<n}(i+1)$.
\end{definition}

\begin{theorem}[Fakultätsregeln]
\label{b10guide:thm:factorial-laws}
Für jedes $n\in\mathbb N$ ist $n!\ge1$.
Für $m\le n$ gelten
\[
 n!=m!\prod_{i<n-m}(m+i+1),\qquad m!\le n!.
\]
Für $1\le m<n$ gilt sogar $m!<n!$; die beiden Anfangswerte
sind $0!=1!=1$.
\end{theorem}

\begin{definition}[Binomialkoeffizienten durch die Pascal-Rekursion]
\label{b10guide:def:binomial}
Für $n,k\in\mathbb N$ sind die natürlichen Zahlen $\binom nk$
eindeutig durch
\[
 \binom00=1,\qquad \binom0{k+1}=0,\qquad
 \binom{n+1}0=1,
\]
und
\[
 \binom{n+1}{k+1}=\binom nk+\binom n{k+1}
\]
bestimmt.
\end{definition}

\begin{theorem}[Grundformeln der Binomialkoeffizienten]
\label{b10guide:thm:binomial-laws}
Es gelten
\[
 \binom n0=\binom nn=1,\qquad
 \binom nk=0\ \Longleftrightarrow\ n<k,
 \qquad \binom n1=n.
\]
Für $k\le n$ gelten
\[
 \binom nk=\binom n{n-k},\qquad
 k!(n-k)!\binom nk=n!.
\]
Die letzte Gleichung charakterisiert $\binom nk$ eindeutig als
natürliche Zahl, ohne einen anderen Zahlbereich einzuführen.
\end{theorem}

\begin{theorem}[Binomischer Lehrsatz]
\label{b10guide:thm:binomial-theorem}
Für $a,b,n\in\mathbb N$ gilt
\[
 (a+b)^n=\sum_{k=0}^{n}\binom nk a^k b^{n-k}.
\]
Die Konvention $0^0=1$ schließt dabei alle Randfälle ein.
\end{theorem}

\begin{theorem}[Zeilensummen und Summen entlang einer Diagonale]
\label{b10guide:thm:binomial-sums}
Für $n,r\in\mathbb N$ gelten
\[
 \sum_{k=0}^{n}\binom nk=2^n,\qquad
 \sum_{i=0}^{n}\binom{i+r}r=\binom{n+r+1}{r+1}.
\]
\end{theorem}


\begin{theorem}[Vandermondesche Identität]
\label{b10guide:thm:vandermonde}
Für $m,n,r\in\mathbb N$ gilt
\[
 \sum_{k=0}^{r}\binom mk\binom n{r-k}=\binom{m+n}r.
\]
Insbesondere gilt für jedes $n\in\mathbb N$
\[
 \sum_{k=0}^{n}\binom nk^2=\binom{2n}n.
\]
Die Definition der Binomialkoeffizienten erfaßt dabei auch die Fälle,
in denen die untere Zahl die obere übersteigt \cite{natvandermonde}.
\end{theorem}

\section{Potenzdifferenzen und natürliche Wurzelwerte}

Die folgenden Aussagen ergänzen die Potenzgesetze um eine
Faktorisierungsformel und eine durch natürliche Potenzschranken
bestimmte Zahl. Vergleichsquellen sind \cite{natgeomsum,natroot}.

\begin{theorem}[Faktorisierung einer Potenzdifferenz]
\label{b10guide:thm:power-difference}
Für $a,b,n\in\mathbb N$ mit $b\le a$ und $n\ge1$ gilt
\[
 a^n-b^n=(a-b)\sum_{i<n}a^{n-1-i}b^i.
\]
Alle Differenzen in dieser Formel sind unter den genannten
Voraussetzungen in $\mathbb N$ definiert.
\end{theorem}

\begin{theorem}[Eindeutige natürliche Wurzelschranke]
\label{b10guide:thm:natural-root-existence}
Zu jeder Zahl $n\in\mathbb N$ und jedem $r\in\mathbb N_+$
gibt es genau eine Zahl $q\in\mathbb N$ mit
\[
 q^r\le n<(q+1)^r.
\]
\end{theorem}

\begin{definition}[Natürlicher Wurzelwert]
\label{b10guide:def:natural-root}
Für $r\in\mathbb N_+$ und $n\in\mathbb N$ bezeichnet
$\operatorname{wurzel}_r(n)$ die eindeutig bestimmte Zahl $q$
aus Satz~\ref{b10guide:thm:natural-root-existence}.
Sie ist durch die beiden Potenzschranken $q^r\le n<(q+1)^r$
definiert. Insbesondere gelten
\[
 \operatorname{wurzel}_1(n)=n,\qquad
 \operatorname{wurzel}_r(0)=0,\qquad
 \operatorname{wurzel}_r(1)=1.
\]
\end{definition}

\begin{theorem}[Vergleich und Erkennung genauer Potenzen]
\label{b10guide:thm:natural-root-characterization}
Für $r\in\mathbb N_+$ und $a,n,q\in\mathbb N$ gilt
\[
\begin{aligned}
 q\le\operatorname{wurzel}_r(n)&\ \Longleftrightarrow\ q^r\le n,\\
 \operatorname{wurzel}_r(n)<q&\ \Longleftrightarrow\ n<q^r,\\
 \operatorname{wurzel}_r(a^r)&=a.
\end{aligned}
\]
Eine natürliche Zahl $n$ ist genau dann eine $r$-te Potenz
einer natürlichen Zahl, wenn
\[
 \bigl(\operatorname{wurzel}_r(n)\bigr)^r=n.
\]
In diesem Fall ist ihre natürliche Basis eindeutig.
\end{theorem}

